% !TEX root = ../LosningsforslagTRES0312.tex
\chapter{Dynamikk}
\label{LF:Chapter:dynamikk}
Dette kapittelet inneholder løsningsforslag til oppgavene i Kapittel~\ref{komp-Chapter:dynamikk}, \textit{Dynamikk}, i kompendiet.

\oppgaver{Geværkule og bil}{k4oppg1}{%
%
En typisk geværkule har masse $m_k = \SI{10}{\gram}$ og forlater løpet med fart $v_k = \SI{900}{\metre\per\second}$.
En bil har masse $m_b = \SI{1500}{\kilo\gram}$.
\begin{enumerate}[a)]
    \item Finn bevegelsesmengden til geværkula.
    \item Hvor fort må bilen kjøre for å ha nøyaktig samme bevegelsesmengde som kula?
          Gi svaret i både \si{\metre\per\second} og \si{\kilo\metre\per\hour}.
\end{enumerate}
%
}

\svar[title={Løsningsforslag}]{%
\noindent
\textbf{a)} Bevegelsesmengden er $p=m\cdot v$ (se \textit{Bevegelsesmengde} i Kapittel~\ref{komp-Chapter:dynamikk}). Vi gjør om massen til kg, $m_k=\SI{10}{\gram}=\SI{0.010}{\kilo\gram}$:
\[
    p = m_k v_k = \SI{0.010}{\kilo\gram}\cdot\SI{900}{\metre\per\second} = \doubleunderline{\SI{9.0}{\kilo\gram\metre\per\second}}
\]

\smallskip\noindent
\textbf{b)} Bilen skal ha samme bevegelsesmengde, $p = m_b v_b$. Vi løser for $v_b$:
\[
    v_b = \frac{p}{m_b} = \frac{\SI{9.0}{\kilo\gram\metre\per\second}}{\SI{1500}{\kilo\gram}} = \doubleunderline{\SI{0.0060}{\metre\per\second}} = \doubleunderline{\SI{0.0216}{\kilo\metre\per\hour}}
\]
Bilen trenger altså bare å kjøre med en helt ubetydelig fart (under 2.2 cm/s!) for å ha samme bevegelsesmengde som den raske, men svært lette, geværkula. Dette illustrerer hvor mye massen betyr for bevegelsesmengden.
}

\oppgaver{Finn fjærkonstanten}{k4oppg2}{%
%
Du fester et lodd til en vertikal fjær og måler at fjæra strekkes \SI{12}{\centi\metre} når loddet har masse $m = \SI{250}{\gram}$.
\begin{enumerate}[a)]
    \item Hva er fjærkonstanten $k$? (Hint: når fjæra er i ro, er fjærkraften lik tyngdekraften på loddet.)
    \item Hvor mye strekkes fjæra om du bytter til et lodd med masse $m = \SI{400}{\gram}$?
\end{enumerate}
%
}

\svar[title={Løsningsforslag}]{%
\noindent
\textbf{a)} Loddet henger i ro, så fjærkraften er like stor som tyngdekraften på loddet (se \textit{Hookes lov} i Kapittel~\ref{komp-Chapter:dynamikk}): $F=kx=m\mathrm{g}$. Vi løser for $k$:
\[
    k = \frac{m\mathrm{g}}{x} = \frac{\SI{0.250}{\kilo\gram}\cdot\SI{9.81}{\metre\per\second\squared}}{\SI{0.12}{\metre}} = \doubleunderline{\SI{20.4}{\newton\per\metre}}
\]

\smallskip\noindent
\textbf{b)} Med det nye loddet, $m=\SI{0.400}{\kilo\gram}$, løser vi $kx=m\mathrm{g}$ for $x$:
\[
    x = \frac{m\mathrm{g}}{k} = \frac{\SI{0.400}{\kilo\gram}\cdot\SI{9.81}{\metre\per\second\squared}}{\SI{20.4}{\newton\per\metre}} = \doubleunderline{\SI{0.192}{\metre}} = \doubleunderline{\SI{19.2}{\centi\metre}}
\]
}

\oppgaver{Tyngdekraft på jord og måne}{k4oppg3}{%
%
Ola har masse $m = \SI{75}{\kilo\gram}$.
\begin{enumerate}[a)]
    \item Finn tyngdekraften på Ola på jordoverflaten ($\mathrm{g} = \SI{9.81}{\metre\per\second\squared}$).
    \item Finn tyngdekraften på Ola på månens overflate ($\mathrm{g}_\text{måne} = \SI{1.62}{\metre\per\second\squared}$).
    \item Hva er Olas masse på månen? Forklar svaret.
\end{enumerate}
%
}

\svar[title={Løsningsforslag}]{%
\noindent
Vi bruker $G=m\mathrm{g}$ (se \textit{Tyngdekraft} i Kapittel~\ref{komp-Chapter:dynamikk}).

\smallskip\noindent
\textbf{a)} På jorda:
\[
    G = m\mathrm{g} = \SI{75}{\kilo\gram}\cdot\SI{9.81}{\metre\per\second\squared} = \doubleunderline{\SI{736}{\newton}}
\]

\smallskip\noindent
\textbf{b)} På månen:
\[
    G_\text{måne} = m\,\mathrm{g}_\text{måne} = \SI{75}{\kilo\gram}\cdot\SI{1.62}{\metre\per\second\squared} = \doubleunderline{\SI{122}{\newton}}
\]

\smallskip\noindent
\textbf{c)} Olas masse er fortsatt $\doubleunderline{m=\SI{75}{\kilo\gram}}$ -- helt uendret. Massen er et mål på hvor mye \textit{treghet} (stoffmengde) Ola har, og det avhenger ikke av hvor han befinner seg. Det er \textit{tyngdekraften} $G$ som endrer seg mellom jorda og månen (fordi $\mathrm{g}$ er forskjellig), ikke massen $m$.
}

\oppgaver{Fritt fall}{k4oppg4}{%
%
En hammer og en fjær slippes samtidig fra samme høyde.
\begin{enumerate}[a)]
    \item Hva tror du skjer — hvilket objekt når bakken først?
    \item Hva skjer om vi pumper vekk all lufta i rommet? Bruk det du har lært om tyngdekraft til å forklare svaret.
\end{enumerate}
%
}

\svar[title={Løsningsforslag}]{%
\noindent
\textbf{a)} De fleste vil tro at hammeren treffer bakken først. Dette stemmer faktisk i virkeligheten (med luft til stede): fjæra har stort overflateareal i forhold til massen sin, så luftmotstanden bremser fjæra mye mer enn den bremser hammeren. Hammeren faller dermed tilnærmet i fritt fall, mens fjæra daler sakte ned.

\smallskip\noindent
\textbf{b)} Dersom vi pumper vekk all lufta (fullstendig vakuum), er luftmotstanden borte, og begge objektene er i \textit{fritt fall} -- de påvirkes kun av tyngdekraften. Fra Newtons 2. lov (se \textit{Tyngdekraft} i Kapittel~\ref{komp-Chapter:dynamikk}):
\[
    \sum F = ma \quad\Rightarrow\quad m\mathrm{g} = ma \quad\Rightarrow\quad a = \mathrm{g},
\]
og massen $m$ forsvinner fullstendig fra likningen. Akselerasjonen er dermed $\mathrm{g}$ for \textit{alle} legemer, uavhengig av masse og form. I vakuum vil derfor hammeren og fjæra \doubleunderline{treffe bakken samtidig} -- noe som faktisk ble demonstrert på Månen av astronaut David Scott under Apollo 15-ferden, nettopp fordi Månen ikke har noen atmosfære.
}

\oppgaver{Normalkraft i heis — ulike situasjoner}{k4oppg5}{%
%
Ola har masse $m = \SI{80}{\kilo\gram}$ og er i en heis. Finn normalkraften $N$ i følgende tilfeller:
\begin{enumerate}[a)]
    \item Heisen bremser ned mens den beveger seg oppover, slik at $a = \SI{-3.0}{\metre\per\second\squared}$.
    \item Heisen er i fritt fall (tau røk!). Hva er $N$ da, og hva ville Ola oppleve?
\end{enumerate}
%
}

\svar[title={Løsningsforslag}]{%
\noindent
Vi velger oppover som positiv retning og bruker Newtons 2. lov i vertikal retning, akkurat som i eksempelet \textit{Normalkraft i akselererende heis} i Kapittel~\ref{komp-Chapter:dynamikk}: $\sum F_y = N - G = ma \implies N = m(\mathrm{g}+a)$.

\smallskip\noindent
\textbf{a)} Her er $a=\SI{-3.0}{\metre\per\second\squared}$ (akselerasjonen peker nedover, siden heisen bremser mens den beveger seg oppover):
\[
    N = m(\mathrm{g}+a) = \SI{80}{\kilo\gram}\cdot(\SI{9.81}{\metre\per\second\squared}-\SI{3.0}{\metre\per\second\squared}) = \doubleunderline{\SI{545}{\newton}}
\]
Dette er mindre enn tyngdekraften ($G=\SI{785}{\newton}$) -- Ola føler seg lettere enn vanlig.

\smallskip\noindent
\textbf{b)} I fritt fall er $a=-\mathrm{g}$:
\[
    N = m(\mathrm{g}-\mathrm{g}) = \doubleunderline{\SI{0}{\newton}}
\]
Normalkraften er null -- Ola ville oppleve \doubleunderline{vektløshet}. Dette er nøyaktig samme situasjon som å simulere vektløshet i fly ved fritt fall (se \textit{Fritt fall som eksempel} i Kapittel~\ref{komp-Chapter:kinematikk}).
}

\oppgaver{Flytte et skap}{k4oppg6}{%
%
Et skap har masse $m = \SI{80}{\kilo\gram}$ og står på et tregulv med $\mu_s = 0.35$.
\begin{enumerate}[a)]
    \item Hva er den minste horisontale kraften som trengs for å sette skapet i bevegelse?
    \item To personer skyver med en samlet kraft på $F = \SI{260}{\newton}$. Begynner skapet å bevege seg? Forklar.
\end{enumerate}
%
}

\svar[title={Løsningsforslag}]{%
\noindent
Normalkraften er $N=m\mathrm{g}=\SI{80}{\kilo\gram}\cdot\SI{9.81}{\metre\per\second\squared}=\SI{784.8}{\newton}$.

\smallskip\noindent
\textbf{a)} Den minste kraften som trengs for å sette skapet i bevegelse må akkurat overvinne maksimal hvilefriksjon (se \textit{Friksjonskrefter} i Kapittel~\ref{komp-Chapter:dynamikk}), $R_{s,\text{maks}}=\mu_s N$:
\[
    F_\text{min} = \mu_s N = 0.35\cdot\SI{784.8}{\newton} = \doubleunderline{\SI{275}{\newton}}
\]

\smallskip\noindent
\textbf{b)} De to personene skyver med $F=\SI{260}{\newton}$, som er \textit{mindre} enn $R_{s,\text{maks}}=\SI{275}{\newton}$. Hvilefriksjonen tilpasser seg og motvirker hele den påførte kraften ($R_s=F=\SI{260}{\newton} < R_{s,\text{maks}}$), så skapet \doubleunderline{begynner ikke å bevege seg}.
}

\oppgaver{Hockeypuck på is}{k4oppg7}{%
%
En hockeypuck har masse $m = \SI{170}{\gram}$ og glir på is med $\mu_k = 0.030$.
\begin{enumerate}[a)]
    \item Hva er friksjonskraften på pucken?
    \item Hva er puckens akselerasjon?
    \item Pucken skyves avgårde med startfart $v_0 = \SI{15}{\metre\per\second}$. Bruk kinematikk til å beregne hvor lang strekning den glir før den stopper.
\end{enumerate}
%
}

\svar[title={Løsningsforslag}]{%
\noindent
\textbf{a)} Normalkraften er $N=m\mathrm{g}$, med $m=\SI{170}{\gram}=\SI{0.170}{\kilo\gram}$. Friksjonskraften er (se \textit{Glidefriksjon} i Kapittel~\ref{komp-Chapter:dynamikk}):
\[
    R = \mu_k N = \mu_k m\mathrm{g} = 0.03\cdot\SI{0.170}{\kilo\gram}\cdot\SI{9.81}{\metre\per\second\squared} = \doubleunderline{\SI{0.0500}{\newton}}
\]

\smallskip\noindent
\textbf{b)} Friksjonen er eneste horisontale kraft, så akselerasjonen er (uavhengig av massen):
\[
    a = -\mu_k \mathrm{g} = -0.03\cdot\SI{9.81}{\metre\per\second\squared} = \doubleunderline{\SI{-0.294}{\metre\per\second\squared}}
\]

\smallskip\noindent
\textbf{c)} Vi bruker fart-strekning-likningen (Likning 3 i avsnitt~\ref{komp-sec:bevegelseslikninger}), $v^2=v_0^2+2a\,\Delta s$, med $v=0$ (pucken stopper):
\[
    0 = v_0^2 + 2a\,s \quad\Rightarrow\quad s = \frac{-v_0^2}{2a} = \frac{v_0^2}{2\mu_k \mathrm{g}} = \frac{(\SI{15}{\metre\per\second})^2}{2\cdot 0.03\cdot\SI{9.81}{\metre\per\second\squared}} = \doubleunderline{\SI{382}{\metre}}
\]
Pucken glir svært langt, nettopp fordi friksjonskoeffisienten mellom puck og is er så liten.
}

\oppgaver{Bremselengde for bildekk}{k4oppg8}{%
%
Et kjøretøy bremser hardt på flatt underlag slik at hjulene låser seg og bildekket sklir. Startfarten er $v_0 = \SI{80}{\kilo\metre\per\hour}$.
\begin{enumerate}[a)]
    \item Vis ved hjelp av Newtons 2. lov og kinematikk at bremselengden er gitt ved
    \[
        s = \frac{v_0^2}{2\,\mu_k\,\mathrm{g}}.
    \]
    \item Beregn bremselengden for tørr asfalt ($\mu_k = 0.70$), våt vei ($\mu_k = 0.40$) og is ($\mu_k = 0.10$).
    \item Hva skjer med bremselengden om farten dobles? Forklar med formelen.
\end{enumerate}
%
}

\svar[title={Løsningsforslag}]{%
\noindent
\textbf{a)} Når hjulene låser seg, glir dekket mot veien, og glidefriksjonen er eneste horisontale kraft. Newtons 2. lov i horisontal retning gir ($x$ i bevegelsesretningen, $N=m\mathrm{g}$ fra vertikal likevekt):
\[
    \sum F_x = -R = ma \quad\Rightarrow\quad -\mu_k m\mathrm{g} = ma \quad\Rightarrow\quad a = -\mu_k\mathrm{g}
\]
(dette er samme resultat som i \textit{Glidefriksjon}, Kapittel~\ref{komp-Chapter:dynamikk}). Deretter bruker vi fart-strekning-likningen (Likning 3 i avsnitt~\ref{komp-sec:bevegelseslikninger}), $v^2=v_0^2+2a\,s$, med $v=0$ når bilen har stoppet:
\[
    0 = v_0^2 - 2\mu_k\mathrm{g}\,s \quad\Longrightarrow\quad \doubleunderline{s = \frac{v_0^2}{2\,\mu_k\,\mathrm{g}}} \qquad\checkmark
\]

\smallskip\noindent
\textbf{b)} Vi regner om $v_0=\SI{80}{\kilo\metre\per\hour}=\SI{22.2}{\metre\per\second}$, og setter inn for hver verdi av $\mu_k$:
\[
    \mu_k=0.70:\quad s = \frac{22.2^2}{2\cdot0.70\cdot9.81} = \doubleunderline{\SI{36.0}{\metre}} \quad\text{(tørr asfalt)}
\]
\[
    \mu_k=0.40:\quad s = \frac{22.2^2}{2\cdot0.40\cdot9.81} = \doubleunderline{\SI{62.9}{\metre}} \quad\text{(våt vei)}
\]
\[
    \mu_k=0.10:\quad s = \frac{22.2^2}{2\cdot0.10\cdot9.81} = \doubleunderline{\SI{252}{\metre}} \quad\text{(is)}
\]

\smallskip\noindent
\textbf{c)} Siden $s \propto v_0^2$, vil en dobling av farten ($v_0 \to 2v_0$) gi
\[
    s' = \frac{(2v_0)^2}{2\mu_k\mathrm{g}} = 4\cdot\frac{v_0^2}{2\mu_k\mathrm{g}} = 4s,
\]
altså \doubleunderline{firedobles} bremselengden. Dette er grunnen til at det er så mye farligere å kjøre fort -- bremselengden øker med \textit{kvadratet} av farten.
}

\oppgaver{Kjelke med skrå trekkraft}{k4oppg9}{%
%
En kjelke med masse $m = \SI{15}{\kilo\gram}$ dras langs et flatt, snødekket underlag. Trekkraften $\vec{F} = \SI{50}{\newton}$ er rettet i vinkel $\theta = 30°$ over horisontal (se figur). Den kinetiske friksjonskoeffisienten er $\mu_k = 0.20$.
\begin{enumerate}[a)]
    \item Finn normalkraften $N$. (Husk at $\vec{F}$ har en vertikal komponent som påvirker $N$.)
    \item Finn friksjonskraften $R$.
    \item Finn kjelkens akselerasjon.
\end{enumerate}
%
}

\svar[title={Løsningsforslag}]{%
\noindent
\textbf{1. Figur:} Kjelken påvirkes av fire krefter: tyngdekraften $\vec G$ (ned), normalkraften $\vec N$ (opp), trekkraften $\vec F$ (i vinkel $\theta=30^\circ$ over horisontalen) og friksjonskraften $\vec R$ (mot bevegelsesretningen), akkurat som i figur~\ref{komp-fig:kjelke}.

\smallskip\noindent
\textbf{2. Lov:} Newtons 2. lov, $\sum\vec F = m\vec a$. Friksjonen er glidefriksjon: $R=\mu_k N$.

\smallskip\noindent
\textbf{3. Koordinatsystem:} $x$ horisontalt (i bevegelsesretningen), $y$ vertikalt oppover. Kjelken beveger seg bare langs bakken, så $a_y=0$.

\smallskip\noindent
\textbf{4. Dekomponér:} Trekkraften har komponenter $F_x=F\cos\theta$ og $F_y=F\sin\theta$. I $y$-retning:
\[
    \sum F_y = N + F\sin\theta - m\mathrm{g} = 0 \quad\Rightarrow\quad N = m\mathrm{g} - F\sin\theta
\]
I $x$-retning:
\[
    \sum F_x = F\cos\theta - R = ma
\]

\smallskip\noindent
\textbf{5. Løs:}

\textbf{a)} Normalkraften:
\[
    N = m\mathrm{g} - F\sin\theta = \SI{15}{\kilo\gram}\cdot\SI{9.81}{\metre\per\second\squared} - \SI{50}{\newton}\cdot\sin(30^\circ) = \doubleunderline{\SI{122}{\newton}}
\]
Legg merke til at $N$ er \textit{mindre} enn tyngdekraften, siden trekkraftens vertikale komponent bidrar til å løfte kjelken litt.

\textbf{b)} Friksjonskraften:
\[
    R = \mu_k N = 0.20\cdot\SI{122}{\newton} = \doubleunderline{\SI{24.4}{\newton}}
\]

\textbf{c)} Akselerasjonen, fra $x$-likningen:
\[
    a = \frac{F\cos\theta - R}{m} = \frac{\SI{50}{\newton}\cdot\cos(30^\circ) - \SI{24.4}{\newton}}{\SI{15}{\kilo\gram}} = \doubleunderline{\SI{1.26}{\metre\per\second\squared}}
\]
}

\oppgaver{Buss på glatt vei}{k4oppg10}{%
%
En buss parkerer på en bakke. Friksjonskoeffisientene mellom dekk og vei er $\mu_s = 0.70$ (statisk) og $\mu_k = 0.50$ (kinetisk).
\begin{enumerate}[a)]
    \item Ved hvilken helningsvinkel $\theta_c$ begynner bussen å skli nedover?
    \item Hva blir akselerasjonen til bussen idet den begynner å skli?
\end{enumerate}
%
}

\svar[title={Løsningsforslag}]{%
\noindent
\textbf{a)} Bussen begynner å skli idet tyngdekraftens komponent langs bakken akkurat overgår maksimal hvilefriksjon -- dette er den kritiske vinkelen utledet i \textit{Kritisk vinkel} (avsnittet \textit{Skråplan} i Kapittel~\ref{komp-Chapter:dynamikk}):
\[
    \theta_c = \arctan(\mu_s) = \arctan(0.70) = \doubleunderline{35.0^\circ}
\]

\smallskip\noindent
\textbf{b)} Idet bussen begynner å skli, glir den, så vi bruker \textit{kinetisk} friksjon $\mu_k$ i formelen for akselerasjon ned et skråplan (utledet i samme avsnitt):
\[
    a = \mathrm{g}(\sin\theta_c - \mu_k\cos\theta_c) = \SI{9.81}{\metre\per\second\squared}\cdot\left(\sin(35.0^\circ) - 0.50\cdot\cos(35.0^\circ)\right) = \doubleunderline{\SI{1.61}{\metre\per\second\squared}}
\]
Akselerasjonen er rettet nedover bakken.
}

\oppgaver{Karusell}{k4oppg11}{%
%
Et barn med masse $m = \SI{25}{\kilo\gram}$ sitter $r = \SI{3.0}{\metre}$ fra sentrum av en karusell som snurrer horisontalt. Karusellen snurrer med konstant fart $v = \SI{4.0}{\metre\per\second}$.

Hva er den resulterende sentripetalkraften på barnet?
%
}

\svar[title={Løsningsforslag}]{%
\noindent
Den resulterende kraftsummen på barnet må peke inn mot sentrum og ha størrelse $\sum F = mv^2/r$ (se \textit{Sentripetalkraft} i Kapittel~\ref{komp-Chapter:dynamikk}):
\[
    \sum F = \frac{mv^2}{r} = \frac{\SI{25}{\kilo\gram}\cdot(\SI{4.0}{\metre\per\second})^2}{\SI{3.0}{\metre}} = \doubleunderline{\SI{133}{\newton}}
\]
Denne kraften må komme fra noe som faktisk trekker eller skyver barnet inn mot sentrum -- for eksempel et sikkerhetsbelte eller friksjon mot setet.
}

\oppgaver{To skøyteløpere}{k4oppg12}{%
%
Oda ($m_O = \SI{55}{\kilo\gram}$) og Per ($m_P = \SI{80}{\kilo\gram}$) står i ro på isen og skyver fra hverandre. Etter skyvingen beveger Oda seg med fart $\SI{4.0}{\metre\per\second}$ i én retning.
\begin{enumerate}[a)]
    \item Hva er Pers hastighet etter skyvingen?
    \item Hva er bevegelsesmengden til hvert av dem etter skyvingen? Kommenter svaret.
\end{enumerate}
%
}

\svar[title={Løsningsforslag}]{%
\noindent
Oda og Per utgjør et isolert system (vi ser bort fra friksjon mot isen), så bevegelsesmengden er bevart (se avsnitt~\ref{komp-sec:bevaring}), akkurat som i eksempelet \textit{Rekyl fra gevær}. Før skyvingen er alt i ro, $p_{tot,\,\text{før}}=0$. Vi velger Odas bevegelsesretning som positiv.

\smallskip\noindent
\textbf{a)} Bevaringslikningen gir
\[
    0 = m_O u_O + m_P u_P \quad\Rightarrow\quad u_P = -\frac{m_O u_O}{m_P} = -\frac{\SI{55}{\kilo\gram}\cdot\SI{4.0}{\metre\per\second}}{\SI{80}{\kilo\gram}} = \doubleunderline{\SI{-2.75}{\metre\per\second}}
\]
Per beveger seg altså med $\SI{2.75}{\metre\per\second}$ i \textit{motsatt} retning av Oda.

\smallskip\noindent
\textbf{b)} Bevegelsesmengden til hver av dem:
\[
    p_O = m_O u_O = \SI{55}{\kilo\gram}\cdot\SI{4.0}{\metre\per\second} = \doubleunderline{\SI{220}{\kilo\gram\metre\per\second}}
\]
\[
    p_P = m_P u_P = \SI{80}{\kilo\gram}\cdot(\SI{-2.75}{\metre\per\second}) = \doubleunderline{\SI{-220}{\kilo\gram\metre\per\second}}
\]
Bevegelsesmengdene er like store og motsatt rettet, slik at summen fortsatt er null -- akkurat som før skyvingen. Dette er nøyaktig det bevaringsloven sier skal skje.
}

\oppgaver{Bilkollisjon}{k4oppg13}{%
%
En bil med masse $m_1 = \SI{1200}{\kilo\gram}$ kjører med fart $\SI{20}{\metre\per\second}$ mot en bil med masse $m_2 = \SI{1000}{\kilo\gram}$ som kjører i \textit{motsatt} retning med fart $\SI{15}{\metre\per\second}$. Bilene henger fast etter kollisjonen.
\begin{enumerate}[a)]
    \item Hva er den totale bevegelsesmengden til systemet \textit{før} kollisjonen?
    \item Finn farten og retningen til bilene etter kollisjonen.
\end{enumerate}
%
}

\svar[title={Løsningsforslag}]{%
\noindent
Vi velger bil 1 sin bevegelsesretning som positiv, slik at $v_1=\SI{20}{\metre\per\second}$ og $v_2=\SI{-15}{\metre\per\second}$ (motsatt retning). Dette er et fullstendig uelastisk støt (bilene henger fast), analogt med eksempelet \textit{Uelastisk støt: vogner som kolliderer} i avsnitt~\ref{komp-sec:bevaring}.

\smallskip\noindent
\textbf{a)} Total bevegelsesmengde før kollisjonen:
\[
    p_{tot,\,\text{før}} = m_1 v_1 + m_2 v_2 = \SI{1200}{\kilo\gram}\cdot\SI{20}{\metre\per\second} + \SI{1000}{\kilo\gram}\cdot(\SI{-15}{\metre\per\second}) = \doubleunderline{\SI{9000}{\kilo\gram\metre\per\second}}
\]

\smallskip\noindent
\textbf{b)} Siden bevegelsesmengden er bevart, og bilene beveger seg som én masse $m_1+m_2$ etter kollisjonen:
\[
    p_{tot,\,\text{før}} = (m_1+m_2)\,u \quad\Rightarrow\quad u = \frac{p_{tot,\,\text{før}}}{m_1+m_2} = \frac{\SI{9000}{\kilo\gram\metre\per\second}}{\SI{2200}{\kilo\gram}} = \doubleunderline{\SI{4.09}{\metre\per\second}}
\]
Farten er positiv, så bilvraket beveger seg videre i \doubleunderline{samme retning som bil 1} (den tyngste og raskeste av de to) etter sammenstøtet.
}

\oppgaver{Kraft og motkraft}{k4oppg14}{%
%
En bok ligger i ro på et bord.
\begin{enumerate}[a)]
    \item Tegn en figur som viser kreftene som virker \emph{på boka}.
    \item Tyngdekraften og normalkraften på boka er like store og motsatt rettet. Er de et kraft--motkraft-par etter Newtons 3. lov? Forklar.
    \item Hva er motkraften til tyngdekraften på boka? Hva er motkraften til normalkraften på boka?
\end{enumerate}
%
}

\svar[title={Løsningsforslag}]{%
\noindent
\textbf{a)} To krefter virker på boka: tyngdekraften $\vec G$ (ned) og normalkraften $\vec N$ fra bordet (opp):
\begin{center}
\begin{tikzpicture}[>=Stealth, scale=1]
    \draw[black!30, thin] (-1.3,0) -- (1.3,0);
    \fill[headCol!20] (-0.7,0) rectangle (0.7,0.35);
    \draw[headCol!65!black, thick] (-0.7,0) rectangle (0.7,0.35);
    \node[font=\small] at (0,0.175) {bok};
    \draw[->, thick, defscol] (0,0.35) -- (0,1.3) node[above, font=\small] {$\vec N$};
    \draw[->, thick, black!65] (0,0) -- (0,-0.95) node[below, font=\small] {$\vec G$};
\end{tikzpicture}
\end{center}

\smallskip\noindent
\textbf{b)} Nei. Selv om $\vec G$ og $\vec N$ er like store og motsatt rettet, er de \doubleunderline{ikke} et kraft-motkraft-par. Grunnen er at et kraft-motkraft-par (Newtons 3. lov) alltid virker på \textit{to forskjellige} legemer, mens $\vec G$ og $\vec N$ begge virker på \textit{samme} legeme -- boka (se \textit{Normalkraft}, tankeboksen i Kapittel~\ref{komp-Chapter:dynamikk}).

\smallskip\noindent
\textbf{c)} Motkraften til tyngdekraften $\vec G$ (jordas tiltrekning på boka) er boka sin tiltrekning på jorda -- en kraft med samme størrelse $G$, rettet oppover, som virker \textit{på jorda}. Motkraften til normalkraften $\vec N$ (bordets kraft på boka) er boka sin kraft \textit{på bordet} -- en kraft med samme størrelse $N$, rettet nedover, som virker \textit{på bordet} (boka presser ned på bordet med samme kraft som bordet presser opp på boka).
}

\oppgaver{Kloss på skråplan med friksjon}{k4oppg15}{%
%
En kloss med masse $m = \SI{5.0}{\kilo\gram}$ glir nedover et skråplan med helningsvinkel $\theta = 30^\circ$. Den kinetiske friksjonskoeffisienten er $\mu_k = 0.20$.
\begin{enumerate}[a)]
    \item Tegn kreftene på klossen og dekomponer tyngdekraften langs og vinkelrett på skråplanet.
    \item Finn normalkraften $N$.
    \item Finn friksjonskraften $R$.
    \item Finn akselerasjonen til klossen.
\end{enumerate}
%
}

\svar[title={Løsningsforslag}]{%
\noindent
\textbf{1. Figur:} Kreftene på klossen er akkurat som i figur~\ref{komp-fig:skråplankrefter1}: normalkraften $\vec N$ (vinkelrett ut fra planet), tyngdekraften $\vec G=m\vec{\mathrm g}$ (rett ned) og friksjonskraften $\vec R$ (langs planet, oppover siden klossen glir nedover).

\smallskip\noindent
\textbf{a)} Vi dekomponerer tyngdekraften langs planet og vinkelrett på planet:
\[
    G_x = m\mathrm{g}\sin\theta \quad\text{(langs planet, nedover)}, \qquad G_y = m\mathrm{g}\cos\theta \quad\text{(vinkelrett på planet)}
\]

\smallskip\noindent
\textbf{2. Lov:} Newtons 2. lov, $\sum\vec F=m\vec a$. Glidefriksjon: $R=\mu_k N$.

\smallskip\noindent
\textbf{3. Koordinatsystem:} $x$ langs planet (positiv nedover), $y$ vinkelrett ut fra planet. Klossen beveger seg bare langs $x$, så $a_y=0$.

\smallskip\noindent
\textbf{4. Dekomponér:}
\[
    \sum F_y = N - m\mathrm{g}\cos\theta = 0, \qquad \sum F_x = m\mathrm{g}\sin\theta - R = ma
\]

\smallskip\noindent
\textbf{5. Løs:}

\textbf{b)} Normalkraften:
\[
    N = m\mathrm{g}\cos\theta = \SI{5.0}{\kilo\gram}\cdot\SI{9.81}{\metre\per\second\squared}\cdot\cos(30^\circ) = \doubleunderline{\SI{42.5}{\newton}}
\]

\textbf{c)} Friksjonskraften:
\[
    R = \mu_k N = 0.20\cdot\SI{42.5}{\newton} = \doubleunderline{\SI{8.50}{\newton}}
\]

\textbf{d)} Akselerasjonen (se også \textit{Utledning av akselerasjon ned et skråplan med friksjon} i Kapittel~\ref{komp-Chapter:dynamikk}):
\[
    a = \mathrm{g}(\sin\theta - \mu_k\cos\theta) = \SI{9.81}{\metre\per\second\squared}\cdot\left(\sin(30^\circ) - 0.20\cos(30^\circ)\right) = \doubleunderline{\SI{3.21}{\metre\per\second\squared}}
\]
rettet nedover langs planet.
}

\oppgaver{Bil over en bakketopp}{k4oppg16}{%
%
En bil med masse $m = \SI{1000}{\kilo\gram}$ kjører over en bakketopp. Bakketoppen kan tilnærmes med en sirkelbue med radius $r = \SI{40}{\metre}$. Bilen har fart $v = \SI{15}{\metre\per\second}$ på toppen.
\begin{enumerate}[a)]
    \item Tegn kreftene på bilen på toppen av bakken. Hvilken vei peker kraftsummen?
    \item Finn normalkraften fra veien på bilen.
    \item Hvor fort kan bilen kjøre over toppen før den mister kontakten med veien?
\end{enumerate}
%
}

\svar[title={Løsningsforslag}]{%
\noindent
\textbf{1. Figur:} Bilen befinner seg på \textit{utsiden} av den sirkulære bakketoppen (motsatt av vogna på \textit{innsiden} av loopen i Kapittel~\ref{komp-Chapter:dynamikk}, figur~\ref{komp-fig:loop_tivoli}). Sentrum av sirkelbuen ligger rett under bilen, så ``inn mot sentrum'' betyr her \textit{nedover}. To krefter virker: tyngdekraften $\vec G=m\vec{\mathrm g}$ (ned, altså inn mot sentrum) og normalkraften $\vec N$ fra veien (opp, altså \textit{ut} fra sentrum -- veien kan bare skyve bilen opp, ikke trekke den ned):
\begin{center}
\begin{tikzpicture}[>=Stealth, scale=1]
    \draw[gray!60, thick] (-1.6,0) arc (200:340:1.7 and 0.55);
    \fill[headCol!20] (-0.45,0.42) rectangle (0.45,0.72);
    \draw[headCol!65!black, thick] (-0.45,0.42) rectangle (0.45,0.72);
    \node[font=\small] at (0,0.57) {$m$};
    \draw[->, thick, defscol] (0,0.72) -- (0,1.55) node[above, font=\small] {$\vec N$};
    \draw[->, thick, black!65] (0,0.42) -- (0,-0.45) node[below, font=\small] {$\vec G$};
    \node[font=\scriptsize, black!55, below] at (0,-0.5) {(mot sentrum)};
\end{tikzpicture}
\end{center}
Siden bilen følger en sirkelbane, må kraftsummen inn mot sentrum (nedover) være $mv^2/r$.

\smallskip\noindent
\textbf{2.--4. Lov, koordinatsystem og dekomponering:} Vi velger positiv retning nedover (inn mot sentrum). Newtons 2. lov gir
\[
    \sum F = m\mathrm{g} - N = \frac{mv^2}{r}
\]
(motsatt fortegn på $N$ sammenliknet med toppen av loopen i figur~\ref{komp-fig:loop_tivoli}, fordi normalkraften her peker \textit{bort} fra sentrum i stedet for \textit{inn} mot det).

\smallskip\noindent
\textbf{5. Løs:}

\textbf{b)} Normalkraften:
\[
    N = m\left(\mathrm{g} - \frac{v^2}{r}\right) = \SI{1000}{\kilo\gram}\cdot\left(\SI{9.81}{\metre\per\second\squared} - \frac{(\SI{15}{\metre\per\second})^2}{\SI{40}{\metre}}\right) = \doubleunderline{\SI{4190}{\newton}}
\]

\textbf{c)} Bilen mister kontakten med veien når $N=0$:
\[
    m\mathrm{g} = \frac{mv_\text{maks}^2}{r} \quad\Rightarrow\quad \doubleunderline{v_\text{maks} = \sqrt{\mathrm{g}r} = \sqrt{\SI{9.81}{\metre\per\second\squared}\cdot\SI{40}{\metre}} = \SI{19.8}{\metre\per\second}}
\]
Kjører bilen fortere enn dette, letter den fra veibanen akkurat på toppen av bakken.
}

\oppgaver{Spark på en fotball}{k4oppg17}{%
%
En fotball med masse $m = \SI{0.45}{\kilo\gram}$ ligger i ro og blir sparket. Etter sparket har ballen fart \SI{25}{\metre\per\second}. Foten er i kontakt med ballen i $\Delta t = \SI{8.0}{\milli\second}$.
\begin{enumerate}[a)]
    \item Finn endringen i bevegelsesmengden til ballen.
    \item Finn den gjennomsnittlige kraften fra foten på ballen.
    \item Hvor stor er denne kraften sammenliknet med tyngdekraften på ballen?
\end{enumerate}
%
}

\svar[title={Løsningsforslag}]{%
\noindent
\textbf{a)} Ballen er i ro før sparket, så endringen i bevegelsesmengde er (se \textit{Impuls} i Kapittel~\ref{komp-Chapter:dynamikk}):
\[
    \Delta p = m v - m\cdot 0 = \SI{0.45}{\kilo\gram}\cdot\SI{25}{\metre\per\second} = \doubleunderline{\SI{11.25}{\kilo\gram\metre\per\second}}
\]

\smallskip\noindent
\textbf{b)} Siden $\Delta p = F\,\Delta t$, finner vi den gjennomsnittlige kraften:
\[
    F = \frac{\Delta p}{\Delta t} = \frac{\SI{11.25}{\kilo\gram\metre\per\second}}{\SI{8.0e-3}{\second}} = \doubleunderline{\SI{1410}{\newton}}
\]

\smallskip\noindent
\textbf{c)} Tyngdekraften på ballen er $G=m\mathrm{g}=\SI{0.45}{\kilo\gram}\cdot\SI{9.81}{\metre\per\second\squared}=\SI{4.41}{\newton}$. Forholdet mellom kraftene er
\[
    \frac{F}{G} = \frac{\SI{1410}{\newton}}{\SI{4.41}{\newton}} \approx \doubleunderline{320}
\]
Kraften fra foten er altså omtrent 320 ganger så stor som tyngdekraften på ballen -- et godt eksempel på hvor store krefter som kan oppstå i et kort støt.
}

\oppgaver{Astronaut i rommet}{k4oppg18}{%
%
En astronaut med masse \SI{90}{\kilo\gram} (inkludert romdrakt) svever i ro \SI{20}{\metre} fra romstasjonen. Hun kaster en verktøykasse med masse \SI{2.0}{\kilo\gram} rett \emph{bort} fra romstasjonen med fart \SI{3.0}{\metre\per\second}.
\begin{enumerate}[a)]
    \item Hvilken fart og retning får astronauten etter kastet?
    \item Hvor lang tid tar det før hun når romstasjonen?
\end{enumerate}
%
}

\svar[title={Løsningsforslag}]{%
\noindent
Astronauten og verktøykassa utgjør et isolert system (ingen ytre krefter i verdensrommet), så bevegelsesmengden er bevart (se avsnitt~\ref{komp-sec:bevaring}) -- akkurat som i eksempelet \textit{Rekyl fra gevær}. Før kastet er begge i ro, $p_{tot,\,\text{før}}=0$. Vi velger retningen bort fra romstasjonen (kasteretningen) som positiv.

\smallskip\noindent
\textbf{a)} Bevaringslikningen gir
\[
    0 = M\,U + m_\text{kasse}\,v_\text{kasse} \quad\Rightarrow\quad U = -\frac{m_\text{kasse}\,v_\text{kasse}}{M} = -\frac{\SI{2.0}{\kilo\gram}\cdot\SI{3.0}{\metre\per\second}}{\SI{90}{\kilo\gram}} = \doubleunderline{\SI{-0.0667}{\metre\per\second}}
\]
Astronauten får altså en fart på $\SI{0.0667}{\metre\per\second}$ i \textit{motsatt} retning av kasteretningen, altså tilbake mot romstasjonen.

\smallskip\noindent
\textbf{b)} Hun beveger seg med konstant fart $|U|$ tilbake mot stasjonen, en avstand $d=\SI{20}{\metre}$ unna:
\[
    t = \frac{d}{|U|} = \frac{\SI{20}{\metre}}{\SI{0.0667}{\metre\per\second}} = \doubleunderline{\SI{300}{\second}} = \SI{5.0}{\minute}
\]
}
